\uuid{Gi7p}
\exo7id{1557}
\titre{exo7 1557}
\auteur{ortiz}
\organisation{exo7}
\datecreate{1999-04-01}
\isIndication{false}
\isCorrection{false}
\chapitre{Endomorphisme particulier}
\sousChapitre{Endomorphisme orthogonal}
\cpge{mpsi-mp2i.applications-lineaires.b, pcsi.applications-lineaires.b, ptsi.applications-lineaires.b, mpsi-mp2i.matrices-applications-lineaires.a, pcsi.matrices-applications-lineaires.a, ptsi.matrices-applications-lineaires.a, mp-mpi.endomorphismes-euclidien.b, pc.endomorphismes-euclidiens.b, psi.endomorphismes-euclidien.b, pt.isometries-euclidien.b, mp-mpi.endomorphismes-euclidien.c, pc.endomorphismes-euclidiens.a, psi.endomorphismes-euclidien.a, pt.isometries-euclidien.a, mp-mpi.endomorphismes-euclidien.e, psi.endomorphismes-euclidien.e, pt.isometries-euclidien.d, bcpst2.valeurs-propres-vecteurs-propres.a, bcpst2.valeurs-propres-vecteurs-propres.b, bcpst2.produit-scalaire.c}
\module{Algèbre}
\niveau{L2}
\difficulte{}

\contenu{
\question{
Que peut-on dire d'une matrice carr\'ee r\'eelle
\`a la fois sym\'etrique et orthogonale ?
D\'eterminer la nature et les \'el\'ements
caract\'eristiques de l'endomorphisme de l'espace
vectoriel euclidien canonique $\Rr^3$ de matrice $
A=\frac {1}{7}\left(\begin{smallmatrix}
-2&6&-3\\
6&3&2\\
-3&2&6
\end{smallmatrix}\right)$ dans la base canonique de $\Rr^3.$
}
}
